% Options for packages loaded elsewhere
% Options for packages loaded elsewhere
\PassOptionsToPackage{unicode}{hyperref}
\PassOptionsToPackage{hyphens}{url}
%
\documentclass[
  letterpaper,
  11pt,
  english,
  onehalfspacing,
  headsepline]{MastersDoctoralThesis}
\usepackage{xcolor}
\usepackage{amsmath,amssymb}
\setcounter{secnumdepth}{5}
\usepackage{iftex}
\ifPDFTeX
  \usepackage[T1]{fontenc}
  \usepackage[utf8]{inputenc}
  \usepackage{textcomp} % provide euro and other symbols
\else % if luatex or xetex
  \usepackage{unicode-math} % this also loads fontspec
  \defaultfontfeatures{Scale=MatchLowercase}
  \defaultfontfeatures[\rmfamily]{Ligatures=TeX,Scale=1}
\fi
\usepackage{lmodern}
\ifPDFTeX\else
  % xetex/luatex font selection
\fi
% Use upquote if available, for straight quotes in verbatim environments
\IfFileExists{upquote.sty}{\usepackage{upquote}}{}
\IfFileExists{microtype.sty}{% use microtype if available
  \usepackage[]{microtype}
  \UseMicrotypeSet[protrusion]{basicmath} % disable protrusion for tt fonts
}{}
\makeatletter
\@ifundefined{KOMAClassName}{% if non-KOMA class
  \IfFileExists{parskip.sty}{%
    \usepackage{parskip}
  }{% else
    \setlength{\parindent}{0pt}
    \setlength{\parskip}{6pt plus 2pt minus 1pt}}
}{% if KOMA class
  \KOMAoptions{parskip=half}}
\makeatother
% Make \paragraph and \subparagraph free-standing
\makeatletter
\ifx\paragraph\undefined\else
  \let\oldparagraph\paragraph
  \renewcommand{\paragraph}{
    \@ifstar
      \xxxParagraphStar
      \xxxParagraphNoStar
  }
  \newcommand{\xxxParagraphStar}[1]{\oldparagraph*{#1}\mbox{}}
  \newcommand{\xxxParagraphNoStar}[1]{\oldparagraph{#1}\mbox{}}
\fi
\ifx\subparagraph\undefined\else
  \let\oldsubparagraph\subparagraph
  \renewcommand{\subparagraph}{
    \@ifstar
      \xxxSubParagraphStar
      \xxxSubParagraphNoStar
  }
  \newcommand{\xxxSubParagraphStar}[1]{\oldsubparagraph*{#1}\mbox{}}
  \newcommand{\xxxSubParagraphNoStar}[1]{\oldsubparagraph{#1}\mbox{}}
\fi
\makeatother


\usepackage{longtable,booktabs,array}
\usepackage{calc} % for calculating minipage widths
% Correct order of tables after \paragraph or \subparagraph
\usepackage{etoolbox}
\makeatletter
\patchcmd\longtable{\par}{\if@noskipsec\mbox{}\fi\par}{}{}
\makeatother
% Allow footnotes in longtable head/foot
\IfFileExists{footnotehyper.sty}{\usepackage{footnotehyper}}{\usepackage{footnote}}
\makesavenoteenv{longtable}
\usepackage{graphicx}
\makeatletter
\newsavebox\pandoc@box
\newcommand*\pandocbounded[1]{% scales image to fit in text height/width
  \sbox\pandoc@box{#1}%
  \Gscale@div\@tempa{\textheight}{\dimexpr\ht\pandoc@box+\dp\pandoc@box\relax}%
  \Gscale@div\@tempb{\linewidth}{\wd\pandoc@box}%
  \ifdim\@tempb\p@<\@tempa\p@\let\@tempa\@tempb\fi% select the smaller of both
  \ifdim\@tempa\p@<\p@\scalebox{\@tempa}{\usebox\pandoc@box}%
  \else\usebox{\pandoc@box}%
  \fi%
}
% Set default figure placement to htbp
\def\fps@figure{htbp}
\makeatother


% definitions for citeproc citations
\NewDocumentCommand\citeproctext{}{}
\NewDocumentCommand\citeproc{mm}{%
  \begingroup\def\citeproctext{#2}\cite{#1}\endgroup}
\makeatletter
 % allow citations to break across lines
 \let\@cite@ofmt\@firstofone
 % avoid brackets around text for \cite:
 \def\@biblabel#1{}
 \def\@cite#1#2{{#1\if@tempswa , #2\fi}}
\makeatother
\newlength{\cslhangindent}
\setlength{\cslhangindent}{1.5em}
\newlength{\csllabelwidth}
\setlength{\csllabelwidth}{3em}
\newenvironment{CSLReferences}[2] % #1 hanging-indent, #2 entry-spacing
 {\begin{list}{}{%
  \setlength{\itemindent}{0pt}
  \setlength{\leftmargin}{0pt}
  \setlength{\parsep}{0pt}
  % turn on hanging indent if param 1 is 1
  \ifodd #1
   \setlength{\leftmargin}{\cslhangindent}
   \setlength{\itemindent}{-1\cslhangindent}
  \fi
  % set entry spacing
  \setlength{\itemsep}{#2\baselineskip}}}
 {\end{list}}
\usepackage{calc}
\newcommand{\CSLBlock}[1]{\hfill\break\parbox[t]{\linewidth}{\strut\ignorespaces#1\strut}}
\newcommand{\CSLLeftMargin}[1]{\parbox[t]{\csllabelwidth}{\strut#1\strut}}
\newcommand{\CSLRightInline}[1]{\parbox[t]{\linewidth - \csllabelwidth}{\strut#1\strut}}
\newcommand{\CSLIndent}[1]{\hspace{\cslhangindent}#1}



\setlength{\emergencystretch}{3em} % prevent overfull lines

\providecommand{\tightlist}{%
  \setlength{\itemsep}{0pt}\setlength{\parskip}{0pt}}



 


% -----------------------------------------------------------------------
%  Extra packages and thesis metadata
%  (MastersDoctoralThesis.cls loads: geometry, hyperref, caption, etc.)
% -----------------------------------------------------------------------

\usepackage[utf8]{inputenc}
\usepackage[T1]{fontenc}
\usepackage{mathpazo}          % Palatino font (matches the template)
\usepackage{csquotes}          % Language-aware quotation marks

% Tables
\usepackage{booktabs}
\usepackage{longtable}
\usepackage{array}
\usepackage{multirow}
\usepackage{float}
\usepackage{pdflscape}

% -----------------------------------------------------------------------
% NOTE: Citations are handled by Quarto's pandoc-citeproc engine using
% the CSL file (csl/nlm-citation-sequence-brackets.csl).
% Do NOT add \usepackage{biblatex} here — it conflicts with citeproc.
% -----------------------------------------------------------------------

% Override class caption: use bold label instead of small-caps (sc)
% The class sets labelfont=sc which renders "Figure" as "FIGURE".
\captionsetup{labelfont=bf}

% -----------------------------------------------------------------------
%  Chapter title format: "1. Introduction" on one line (no "Chapter X")
% -----------------------------------------------------------------------
\makeatletter
\renewcommand{\@makechapterhead}[1]{%
  \abovechapterskip
  {\parindent \z@ \chapteralign \normalfont
    \chapterfont
    \ifnum \c@secnumdepth >\m@ne
      \if@mainmatter
        \thechapter.\enspace
      \fi
    \fi
    \interlinepenalty\@M
    #1\par\nobreak
    \chapterbelowskip
  }
  \thispagestyle{\chapter@p@gestyle}
}
\makeatother

% -----------------------------------------------------------------------
%  Suppress Quarto's auto-generated \maketitle — we have a custom title
%  page in tex/frontmatter.tex instead.
% -----------------------------------------------------------------------
\renewcommand{\maketitle}{}

% -----------------------------------------------------------------------
%  Thesis metadata — used by \ttitle, \authorname, \supname, etc.
% -----------------------------------------------------------------------
\thesistitle{Dairy intake and muscle ageing in older women: a prospective
  analysis from the OsteoLaus cohort}
\supervisor{PD Dr.~med.~Dr.~phil.~Patricia Chocano-Bedoya}
\examiner{}
\degree{Master of Science in Bioinformatics and Computational Biology}
% \author is intentionally omitted here; Quarto injects it from _quarto.yml
% which sets \authorname via the class's \author redefinition.
\addresses{}
\subject{Nutritional Epidemiology}
\keywords{sarcopenia, dairy, muscle, ageing, cohort study, OsteoLaus}
\university{University of Fribourg}
\department{Department of Biology}
\group{}
\faculty{Faculty of Science and Medicine}

\AtBeginDocument{
  \hypersetup{pdftitle=\ttitle}
  \hypersetup{pdfauthor=\authorname}
  \hypersetup{pdfkeywords=\keywordnames}
}
\makeatletter
\@ifpackageloaded{bookmark}{}{\usepackage{bookmark}}
\makeatother
\makeatletter
\@ifpackageloaded{caption}{}{\usepackage{caption}}
\AtBeginDocument{%
\ifdefined\contentsname
  \renewcommand*\contentsname{Table of contents}
\else
  \newcommand\contentsname{Table of contents}
\fi
\ifdefined\listfigurename
  \renewcommand*\listfigurename{List of Figures}
\else
  \newcommand\listfigurename{List of Figures}
\fi
\ifdefined\listtablename
  \renewcommand*\listtablename{List of Tables}
\else
  \newcommand\listtablename{List of Tables}
\fi
\ifdefined\figurename
  \renewcommand*\figurename{Figure}
\else
  \newcommand\figurename{Figure}
\fi
\ifdefined\tablename
  \renewcommand*\tablename{Table}
\else
  \newcommand\tablename{Table}
\fi
}
\@ifpackageloaded{float}{}{\usepackage{float}}
\floatstyle{ruled}
\@ifundefined{c@chapter}{\newfloat{codelisting}{h}{lop}}{\newfloat{codelisting}{h}{lop}[chapter]}
\floatname{codelisting}{Listing}
\newcommand*\listoflistings{\listof{codelisting}{List of Listings}}
\makeatother
\makeatletter
\makeatother
\makeatletter
\@ifpackageloaded{caption}{}{\usepackage{caption}}
\@ifpackageloaded{subcaption}{}{\usepackage{subcaption}}
\makeatother
\usepackage{bookmark}
\IfFileExists{xurl.sty}{\usepackage{xurl}}{} % add URL line breaks if available
\urlstyle{same}
\hypersetup{
  pdftitle={Dairy intake and muscle ageing in older women},
  pdfauthor={Anna Boss},
  hidelinks,
  pdfcreator={LaTeX via pandoc}}


\title{Dairy intake and muscle ageing in older women}
\usepackage{etoolbox}
\makeatletter
\providecommand{\subtitle}[1]{% add subtitle to \maketitle
  \apptocmd{\@title}{\par {\large #1 \par}}{}{}
}
\makeatother
\subtitle{A prospective analysis from the OsteoLaus cohort}
\author{Anna Boss}
\date{2026-06-28}
\begin{document}
\maketitle

% -----------------------------------------------------------------------
%  Front matter — inserted before the first Quarto chapter
% -----------------------------------------------------------------------

\frontmatter
\pagestyle{plain}

% -----------------------------------------------------------------------
%  TITLE PAGE
% -----------------------------------------------------------------------
\begin{titlepage}
\begin{center}

\vspace*{.06\textheight}
{\scshape\LARGE \univname\par}\vspace{1.5cm}
\textsc{\Large Master's Thesis}\\[0.5cm]

\HRule \\[0.4cm]
{\huge \bfseries \ttitle\par}\vspace{0.4cm}
\HRule \\[1.5cm]

\begin{minipage}[t]{0.4\textwidth}
\begin{flushleft} \large
\emph{Author:}\\
\authorname
\end{flushleft}
\end{minipage}
\begin{minipage}[t]{0.4\textwidth}
\begin{flushright} \large
\emph{Supervisor:}\\
\supname
\end{flushright}
\end{minipage}\\[3cm]

\vfill
\degreename\\[0.4cm]
\deptname\\[2cm]

\vfill
{\large \today}
\vfill

\end{center}
\end{titlepage}

% -----------------------------------------------------------------------
%  ABSTRACT
% -----------------------------------------------------------------------
\begin{abstract}
\addchaptertocentry{\abstractname}

% Paste your abstract here (≈ 250 words).
% Use \textbf{Background:}, \textbf{Methods:}, \textbf{Results:},
% \textbf{Conclusions:} if a structured abstract is required.

\end{abstract}

% -----------------------------------------------------------------------
%  ACKNOWLEDGEMENTS
% -----------------------------------------------------------------------
\begin{acknowledgements}
\addchaptertocentry{\acknowledgementname}

% Acknowledge supervisors, co-authors, funding, and personal support here.

\end{acknowledgements}

% -----------------------------------------------------------------------
%  TABLE OF CONTENTS / FIGURES / TABLES
% -----------------------------------------------------------------------
\tableofcontents
\listoffigures
\listoftables

% -----------------------------------------------------------------------
%  ABBREVIATIONS
% -----------------------------------------------------------------------
\begin{abbreviations}{ll}

\textbf{ALM}     & \textbf{A}ppendicular \textbf{L}ean \textbf{M}ass \\
\textbf{ALMI}    & \textbf{A}ppendicular \textbf{L}ean \textbf{M}ass \textbf{I}ndex \\
\textbf{BMI}     & \textbf{B}ody \textbf{M}ass \textbf{I}ndex \\
\textbf{CI}      & \textbf{C}onfidence \textbf{I}nterval \\
\textbf{CoLaus}  & \textbf{Co}hort \textbf{Laus}anne \\
\textbf{DXA}     & \textbf{D}ual-energy \textbf{X}-ray \textbf{A}bsorptiometry \\
\textbf{EWGSOP2} & \textbf{E}uropean \textbf{W}orking \textbf{G}roup on
                   \textbf{S}arcopenia in \textbf{O}lder \textbf{P}eople 2 \\
\textbf{HGS}     & \textbf{H}and\textbf{g}rip \textbf{S}trength \\
\textbf{HR}      & \textbf{H}azard \textbf{R}atio \\
\textbf{IQR}     & \textbf{I}nter\textbf{q}uartile \textbf{R}ange \\
\textbf{LMM}     & \textbf{L}inear \textbf{M}ixed \textbf{M}odel \\
\textbf{MET}     & \textbf{M}etabolic \textbf{E}quivalent of \textbf{T}ask \\
\textbf{MPS}     & \textbf{M}uscle \textbf{P}rotein \textbf{S}ynthesis \\
\textbf{OR}      & \textbf{O}dds \textbf{R}atio \\
\textbf{SD}      & \textbf{S}tandard \textbf{D}eviation \\
\textbf{SE}      & \textbf{S}tandard \textbf{E}rror \\

\end{abbreviations}

% -----------------------------------------------------------------------
%  Begin main matter (switches to arabic page numbering)
% -----------------------------------------------------------------------
\mainmatter
\pagestyle{thesis}


\bookmarksetup{startatroot}

\chapter{}\label{section}

\bookmarksetup{startatroot}

\chapter{Introduction}\label{sec-introduction}

\section{The ageing population and the burden of
sarcopenia}\label{the-ageing-population-and-the-burden-of-sarcopenia}

\textbf{The onset of significant skeletal muscle mass decline in adults
represents a complex temporal process that differs between biological
initiation and clinical manifestation.} While sarcopenia often becomes
clinically significant only in the sixth decade of life, the underlying
degenerative processes begin substantially earlier
{[}\citeproc{ref-westerstahl2025}{1},\citeproc{ref-gustafsson2024}{2}{]}.
A large-scale cross-sectional analysis of 50,332 participants aged
45--84 years reveals continuous decline in fat-free thigh muscle volume
across all age groups {[}\citeproc{ref-anderson2025}{3}{]}. This pattern
is also supported by longitudinal evidence
{[}\citeproc{ref-westerstahl2025}{1}{]}.

\textbf{Age-related muscle loss results from multiple interconnected
physiological and metabolic changes, including neural degeneration,
mitochondrial dysfunction, altered amino acid metabolism, and impaired
muscle protein synthesis.} At the neural level, healthy older adults
experience significant motor unit remodelling, with approximately 40\%
fewer motor units by age 71 {[}\citeproc{ref-piasecki2016}{4}{]}. The
progressive loss of muscle mass is fundamentally driven by the dual
mechanisms of muscle fibre atrophy and hypoplasia, largely resulting
from chronic imbalances between muscle protein synthesis (MPS) and
breakdown (MPB) {[}\citeproc{ref-wilkinson2018}{5}{]}.

\section{Postmenopausal women as a particularly vulnerable
population}\label{postmenopausal-women-as-a-particularly-vulnerable-population}

\section{Dairy as a candidate dietary
exposure}\label{dairy-as-a-candidate-dietary-exposure}

\section{Study objective and
hypothesis}\label{study-objective-and-hypothesis}

\bookmarksetup{startatroot}

\chapter{Methods}\label{sec-methods}

\section{Study design and population}\label{study-design-and-population}

I analyzed longitudinal data from the CoLaus\textbar{}PsyCoLaus cohort
and its nested sub-study, OsteoLaus. CoLaus is a population-based cohort
of men and women aged 35 to 75 years at baseline living in Lausanne,
Switzerland, recruited between June 2003 and May 2006 at a single study
center. Follow-up assessments of the CoLaus cohort were conducted at
four intervals, i.e., 2009-2013, 2014-2018, 2018-2021, and 2022-2026.
OsteoLaus is a sub-cohort nested within CoLaus, comprising women who
participated in the CoLaus baseline examination and were invited to an
additional study focused on bone health. The OsteoLaus study included
\(1\,475\) women aged 50 to 80 years at baseline
{[}\citeproc{ref-shevroja2019}{6}{]}.

Eligibility criteria for the OsteoLaus study included written informed
consent, body mass index (BMI) between 15 and 40 kg/m², and availability
of dual-energy X-ray absorptiometry (DXA) data (obtained in vertebrae
and femur) . In addition, ultrasound bone density measurements in the
heel and information on previous osteoporotic fractures had to be
available. Women with known diseases affecting bone metabolism or
skeletal sites were excluded.

Data were collected approximately every 2.5 years at five visits:
baseline (2010-2012, n = \(1\,475\)), second visit (2012-2015, n =
\(1\,349\)), third visit (2015-2018, n = \(1\,242\)), fourth visit
(2017-2020, n = \(1\,104\)), and fifth visit (2020-2022, n = 944). Loss
to follow-up of participants were mainly due to, death,
institutionalisation, refusal to participate, BMI above 40 kg/m², or
severe psychiatric disease {[}\citeproc{ref-vendrami2025a}{7}{]}.

Further details on recruitment and methodology for CoLaus and OsteoLaus
were described by Firmann et al. {[}\citeproc{ref-firmann2008}{8}{]} and
Shevroja et al. {[}\citeproc{ref-shevroja2019}{6}{]}, respectively.
Additional study information is available at the
CoLaus\textbar{}PsyCoLaus
\href{https://www.colaus-psycolaus.ch/}{website}.
\textbf{?@fig-timeline} shows the study population sizes, visit time
intervals, and the variables collected at each visit.

\section{Dataset matching}\label{sec-matching}

Visits from the CoLaus and OsteoLaus cohorts were aligned using a
predefined visit-matching scheme, with OsteoLaus serving as the backbone
dataset. OsteoLaus baseline was paired with CoLaus follow-up 1 (F1),
OsteoLaus visit 3 (V3) with CoLaus follow-up 2 (F2), and OsteoLaus visit
4 (V4) with CoLaus follow-up 3 (F3), corresponding to study time points
T1, T2, and T3, respectively (\textbf{?@fig-timeline}). OsteoLaus visit
2 (V2) was excluded because it had no corresponding CoLaus assessment
and was not part of the longitudinal analysis framework. Similarly, the
baseline visit of CoLaus was excluded. OsteoLaus visit 5 (V5) was
retained as an OsteoLaus-only assessment and defined as time point T4 to
retain two time points for gait speed. For variables available in both
cohorts, the OsteoLaus value was preferentially retained, with the
corresponding CoLaus value used only when the OsteoLaus value was
missing. Variables collected exclusively at CoLaus visits, such as dairy
intake, were retained unchanged. The final dataset included the time
difference between the matched CoLaus and OsteoLaus examination dates
for each study time point.

\section{Exposure}\label{exposure}

\subsection{Dietary assessment}\label{dietary-assessment}

Dietary intake, including dairy consumption, was assessed in the CoLaus
study from the first follow-up visit onwards using a self-administered
semi-quantitative food frequency questionnaire (FFQ). The FFQ was
developed and validated in Geneva and includes 97 food items covering
intake over the previous 4 weeks
{[}\citeproc{ref-bernstein1995}{9},\citeproc{ref-marques-vidal2018}{10}{]}.
Participants reported consumption frequency and portion size for each
item. Portion sizes were defined as seen in \textbf{?@tbl-ffq} for dairy
products. Daily intake (g/day or mL/day) was calculated by combining
frequency and portion size information. Energy and nutrient values were
derived using the French CIQUAL food composition database
{[}\citeproc{ref-ciqual}{11}{]}.

\begin{longtable}[]{@{}
  >{\raggedright\arraybackslash}p{(\linewidth - 6\tabcolsep) * \real{0.1644}}
  >{\raggedright\arraybackslash}p{(\linewidth - 6\tabcolsep) * \real{0.3973}}
  >{\raggedright\arraybackslash}p{(\linewidth - 6\tabcolsep) * \real{0.0822}}
  >{\raggedright\arraybackslash}p{(\linewidth - 6\tabcolsep) * \real{0.3562}}@{}}
\caption{Food Frequency Questionnaire (FFQ) items and portion sizes for
women.}\tabularnewline
\toprule\noalign{}
\begin{minipage}[b]{\linewidth}\raggedright
FFQ number
\end{minipage} & \begin{minipage}[b]{\linewidth}\raggedright
Food item
\end{minipage} & \begin{minipage}[b]{\linewidth}\raggedright
Unit
\end{minipage} & \begin{minipage}[b]{\linewidth}\raggedright
Portion size (S / E / L)
\end{minipage} \\
\midrule\noalign{}
\endfirsthead
\toprule\noalign{}
\begin{minipage}[b]{\linewidth}\raggedright
FFQ number
\end{minipage} & \begin{minipage}[b]{\linewidth}\raggedright
Food item
\end{minipage} & \begin{minipage}[b]{\linewidth}\raggedright
Unit
\end{minipage} & \begin{minipage}[b]{\linewidth}\raggedright
Portion size (S / E / L)
\end{minipage} \\
\midrule\noalign{}
\endhead
\bottomrule\noalign{}
\endlastfoot
1 & Plain yogurt & g & 90 / 180 / 270 \\
2 & Low-fat yogurt & g & 90 / 180 / 270 \\
3 & Fruit/aroma yogurt & g & 90 / 180 / 270 \\
4 & Cottage cheese (0\% fat) & g & 60 / 125 / 200 \\
5 & Cottage cheese, ricotta & g & 50 / 100 / 150 \\
6 & Feta, mozzarella & g & 20 / 40 / 60 \\
7 & Gruyère, tomme, camembert & g & 30 / 60 / 90 \\
8 & Cheese fondue & g & 175 / 350 / 525 \\
52 & Butter & g & 4 / 7 / 10 \\
53 & Cream (35\% fat) & g & 8 / 15 / 23 \\
63 & Cream tart, cream cake & g & 50 / 100 / 150 \\
68 & Ice cream, sorbet & g & 50 / 100 / 150 \\
71 & Butter (for cooking) & g & 4 / 7 / 10 \\
82 & Milk in coffee (0\% fat) & mL & 15 / 30 / 45 \\
83 & Milk in coffee (non 0\% fat) & mL & 15 / 30 / 45 \\
84 & Coffee creamer & mL & 6 / 12 / 18 \\
85 & Milk as drink (0\% fat) & mL & 75 / 150 / 225 \\
86 & Milk as drink (non 0\% fat) & mL & 75 / 150 / 225 \\
\end{longtable}

FFQ = food frequency questionnaire; S = small, E = medium, L = large
portion.

\subsection{Total dairy intake}\label{total-dairy-intake}

Total dairy intake was defined as the sum of milk, cheese, yogurt,
cream, butter, and dairy-based desserts consumed per day. Intake was
expressed in g/day, assuming a density of 1 g/mL for liquid dairy
products to allow harmonization across units. Ice cream (FFQ68) was
reported together with sorbet. The entire amount of FFQ68 reported was
assumed to represent a dairy-based product and was included in total
dairy intake.

\bookmarksetup{startatroot}

\chapter{Results}\label{sec-results}

\section{Study population}\label{study-population}

\section{Baseline characteristics}\label{sec-baseline}

\section{Dairy intake and handgrip strength}\label{sec-hgs}

\section{Dairy intake and appendicular lean mass index}\label{sec-almi}

\section{Dairy intake and gait speed}\label{sec-gait}

\section{Incident sarcopenia}\label{sec-cox}

\section{Sensitivity analyses}\label{sec-sensitivity}

\bookmarksetup{startatroot}

\chapter{Discussion}\label{sec-discussion}

\section{Principal findings}\label{principal-findings}

\section{Comparison with existing
literature}\label{comparison-with-existing-literature}

\section{Biological mechanisms}\label{biological-mechanisms}

\section{Strengths}\label{sec-strengths}

\section{Limitations}\label{sec-limitations}

\section{Implications for public health}\label{sec-implications}

\bookmarksetup{startatroot}

\chapter{Conclusion}\label{sec-conclusion}

\bookmarksetup{startatroot}

\chapter*{References}\label{references}
\addcontentsline{toc}{chapter}{References}

\markboth{References}{References}

\cleardoublepage
\phantomsection
\addcontentsline{toc}{part}{Appendices}
\appendix

\chapter*{Supplementary Material}\label{sec-supplementary}
\addcontentsline{toc}{chapter}{Supplementary Material}

\markboth{Supplementary Material}{Supplementary Material}

\section*{Directed Acyclic Graph (DAG)}\label{sec-dag}
\addcontentsline{toc}{section}{Directed Acyclic Graph (DAG)}

\markright{Directed Acyclic Graph (DAG)}

\section*{Consort flow diagram}\label{sec-consort}
\addcontentsline{toc}{section}{Consort flow diagram}

\markright{Consort flow diagram}

\begin{figure}

\centering{

\includegraphics[width=0.9\linewidth,height=\textheight,keepaspectratio]{appendices/../images/consort_flow.png}

}

\caption{\label{fig-consort}Consort flow diagram}

\end{figure}%

\section*{Supplementary tables}\label{sec-sup-tables}
\addcontentsline{toc}{section}{Supplementary tables}

\markright{Supplementary tables}

\section*{Supplementary analyses}\label{sec-sup-analyses}
\addcontentsline{toc}{section}{Supplementary analyses}

\markright{Supplementary analyses}

\phantomsection\label{refs}
\begin{CSLReferences}{0}{1}
\bibitem[\citeproctext]{ref-westerstahl2025}
\CSLLeftMargin{1. }%
\CSLRightInline{Maria Westerståhl, Gustav Jörnåker, Eva Jansson, Ulrika
Aasa, Michael Ingre, Kaveh Pourhamidi, Brun Ulfhake, Thomas Gustafsson.
Rise and {Fall} of {Physical Capacity} in a {General Population}: {A}
47-{Year Longitudinal Study}. Journal of Cachexia, Sarcopenia and
Muscle. 2025;16(6):e70134.
doi:\href{https://doi.org/10.1002/jcsm.70134}{10.1002/jcsm.70134}}

\bibitem[\citeproctext]{ref-gustafsson2024}
\CSLLeftMargin{2. }%
\CSLRightInline{Thomas Gustafsson, Brun Ulfhake. Aging {Skeletal
Muscles}: {What Are} the {Mechanisms} of {Age-Related Loss} of
{Strength} and {Muscle Mass}, and {Can We Impede Its Development} and
{Progression}? International Journal of Molecular Sciences. 2024
Oct;25(20).
doi:\href{https://doi.org/10.3390/ijms252010932}{10.3390/ijms252010932}}

\bibitem[\citeproctext]{ref-anderson2025}
\CSLLeftMargin{3. }%
\CSLRightInline{David B. Anderson, Aaron J. Beach, Lingxiao Chen, Henry
J. Feng, Marnee J. McKay, Zachary A. Smith, Kenneth A. Weber, Evert Onno
Wesselink, James M. Elliott. What is normal age-related thigh muscle
composition among 45- to 84-year-old adults from the {UK Biobank} study.
GeroScience. 2025 Feb;47(1):1175--85.
doi:\href{https://doi.org/10.1007/s11357-024-01304-y}{10.1007/s11357-024-01304-y}}

\bibitem[\citeproctext]{ref-piasecki2016}
\CSLLeftMargin{4. }%
\CSLRightInline{Mathew Piasecki, Alex Ireland, David A. Jones, Jamie S.
McPhee. Age-dependent motor unit remodelling in human limb muscles.
Biogerontology. 2016 Jun;17(3):485--96.
doi:\href{https://doi.org/10.1007/s10522-015-9627-3}{10.1007/s10522-015-9627-3}}

\bibitem[\citeproctext]{ref-wilkinson2018}
\CSLLeftMargin{5. }%
\CSLRightInline{D. J. Wilkinson, M. Piasecki, P. J. Atherton. The
age-related loss of skeletal muscle mass and function: {Measurement} and
physiology of muscle fibre atrophy and muscle fibre loss in humans.
Ageing Research Reviews. 2018 Nov;47:123--32.
doi:\href{https://doi.org/10.1016/j.arr.2018.07.005}{10.1016/j.arr.2018.07.005}}

\bibitem[\citeproctext]{ref-shevroja2019}
\CSLLeftMargin{6. }%
\CSLRightInline{Enisa Shevroja, Pedro Marques-Vidal, Bérengère
Aubry-Rozier, Gabriel Hans, Fernando Rivadeneira, Olivier Lamy, Didier
Hans. Cohort {Profile}: {The OsteoLaus} study. International Journal of
Epidemiology. 2019 Aug;48(4):1046--1047g.
doi:\href{https://doi.org/10.1093/ije/dyy276}{10.1093/ije/dyy276}}

\bibitem[\citeproctext]{ref-vendrami2025a}
\CSLLeftMargin{7. }%
\CSLRightInline{Colin Vendrami, Elena Gonzalez Rodriguez, Guillaume
Gatineau, Peter Vollenweider, Pedro Marques-Vidal, Olivier Lamy, Didier
Hans, Enisa Shevroja. Prediction of 10-{Year Fragility Fractures Using
Muscle Health Indicators} in {Postmenopausal Women}: {The OsteoLaus
Cohort}. Journal of Cachexia, Sarcopenia and Muscle. 2025
Jun;16(3):e13837.
doi:\href{https://doi.org/10.1002/jcsm.13837}{10.1002/jcsm.13837} PubMed
PMID: \href{https://www.ncbi.nlm.nih.gov/pubmed/40464199}{40464199};
PubMed Central PMCID:
\href{https://www.ncbi.nlm.nih.gov/pmc/articles/PMC12134769}{PMC12134769}.}

\bibitem[\citeproctext]{ref-firmann2008}
\CSLLeftMargin{8. }%
\CSLRightInline{Mathieu Firmann, Vladimir Mayor, Pedro Marques Vidal,
Murielle Bochud, Alain Pécoud, Daniel Hayoz, Fred Paccaud, Martin
Preisig, Kijoung S. Song, Xin Yuan, Theodore M. Danoff, Heide A.
Stirnadel, Dawn Waterworth, Vincent Mooser, Gérard Waeber, Peter
Vollenweider. The {CoLaus} study: A population-based study to
investigate the epidemiology and genetic determinants of cardiovascular
risk factors and metabolic syndrome. BMC Cardiovascular Disorders. 2008
Mar;8(1):6.
doi:\href{https://doi.org/10.1186/1471-2261-8-6}{10.1186/1471-2261-8-6}}

\bibitem[\citeproctext]{ref-bernstein1995}
\CSLLeftMargin{9. }%
\CSLRightInline{Bernstein, Huot, Morabia. Amélioration des performances
d'un questionnaire alimentaire semi-quantitatif comparé à un rappel des
24 heures. Sante Publique. 1995.}

\bibitem[\citeproctext]{ref-marques-vidal2018}
\CSLLeftMargin{10. }%
\CSLRightInline{Pedro Marques-Vidal, Gérard Waeber, Peter Vollenweider,
Idris Guessous. Socio-demographic and lifestyle determinants of dietary
patterns in {French-speaking Switzerland}, 2009--2012. BMC Public
Health. 2018 Jan;18(1):131.
doi:\href{https://doi.org/10.1186/s12889-018-5045-1}{10.1186/s12889-018-5045-1}}

\bibitem[\citeproctext]{ref-ciqual}
\CSLLeftMargin{11. }%
\CSLRightInline{French Agency for Food, Environmental and Occupational
Health \& Safety. Ciqual. French food composition table.
https://ciqual.anses.fr/.}

\end{CSLReferences}




\end{document}
